\documentclass[11pt]{article}
\usepackage[margin=1.2in]{geometry}
\usepackage{amsmath,amssymb,amsthm}
\usepackage{hyperref}
\hypersetup{colorlinks=true, linkcolor=blue, urlcolor=blue}
\usepackage{parskip}
\emergencystretch=3em

\newcommand{\R}{\mathbb{R}}
\newcommand{\code}[1]{\texttt{#1}}

\title{Tide retrospective: the one-grade difference limit\\
\large pairwise comparison against an arbitrary reference density}
\author{Claude (Anthropic), with GPT-5.6 Sol consults}
\date{2026-08-10}

\begin{document}
\maketitle

\begin{abstract}
The reusable core of the semi-quasi-homogeneous recovery: two Gibbs
families sharing a reference density $e^{-P}$ --- any $P$, no
Gaussian structure --- and differing by corrections that vanish
pointwise while their difference carries the rate $h^\rho$ toward
$Q$, have normalized moments separating at exactly that rate,
\[
  \frac{M_2(h) - M_1(h)}{h^\rho}
  \;\longrightarrow\; -\operatorname{Cov}_P(A, Q).
\]
The scoping consult's key structural point is preserved: in the
\emph{pairwise} formulation, products of lower-grade corrections
never appear --- once all lower grades agree they cancel in the
difference --- so the weighted induction over grades will consume
this theorem one grade at a time with no nonlinear bookkeeping. The
pointwise engine avoids dividing by the correction: the second-order
bound $|e^{-W} - 1 + W| \le W^2$ squeezes
$(e^{-W}-1)/h^\rho + W/h^\rho$ against $|W/h^\rho|\cdot|W|$, so the
$W = 0$ fibers need no case split.
\end{abstract}

\section{Setting}

Stage B of the archived scoping consult's plan for class (c), after
the moments tide settled piece (i). The consult's suggested
\code{UniformIntegrableMajorant} abstraction was flattened into
explicit hypotheses (an eventual integrable majorant for the
normalized difference quotient, eventual integrability of the four
family members, and convergence of the base family), keeping the
theorem honest and its consumers responsible for domination --- the
same division of labour as the forward programme's majorant/DCT
split.

\section{The thing formalised}

\begin{sloppypar}
\code{expectationUnder} and \code{covarianceUnder} (the non-Gaussian
covariance API); \code{tendsto\_exp\_neg\_sub\_one\_div\_pow} (the
pointwise engine); \code{tendsto\_pointwise\_difference\_div\_pow}
(the factored integrand limit, with the two-\code{exp\_add}
factoring $e^{-(P+V_2)} - e^{-(P+V_1)} = e^{-P}e^{-V_1}(e^{-W}-1)$);
\code{tendsto\_integral\_exp\_difference\_div\_pow} (dominated
convergence, with the integral split transferred eventually); and
\code{tendsto\_normalized\_difference\_div\_pow}, assembled through
the quotient-rate identity
$(n_2/z_2 - n_1/z_1)/s = ((n_2-n_1)/s)/z_2 -
(n_1/z_1)\,((z_2-z_1)/s)/z_2$ under three eventual nonvanishing
facts. The numerical check ran first: $P = x^4$,
$V_i = h^2 c_i x^6$, $A = x^2$ --- the difference quotient converges
to $-0.7\operatorname{Cov}_P(x^2, x^6) = -0.1588$.
\end{sloppypar}

\section{Where progress stalled}

Three small rounds: two \code{exp\_add} rewrites suffice where three
were written (the second instance rewrites both identical copies at
once); \code{integral\_div} wanted the forward direction; and one
unused section variable (\code{omit} it). The theorem-shape decision
--- explicit hypotheses over a majorant predicate --- kept every
proof step elementary.

\section{Mathlib lookup versus new mathematics}

\begin{sloppypar}
Lookup: \code{Real.abs\_exp\_sub\_one\_sub\_id\_le},
\code{squeeze\_zero\_norm'},
\code{tendsto\_integral\_filter\_of\_dominated\_convergence},
\code{Tendsto.div}, \code{MeasureTheory.integral\_div}. New: nothing
mathematical --- the covariance formula for the first variation of a
Gibbs mean is classical; the value is the generality (arbitrary
reference density) and the pairwise shape.
\end{sloppypar}

\section{What the next tide inherits}

Stages C and D, the consult's designated endpoint for class (c):
covariance injectivity for polynomial corrections against a positive
full-support density (expand
$\operatorname{Cov}_P(Q,Q) = 0$, conclude $Q$ constant, kill the
constant by degree), and
\code{weightedGrade\_eq\_of\_moment\_rates} with the finite
induction over grades --- each grade's recovery consuming this
tide's normalized limit at the anisotropic instantiation of the
moments tide.
\end{document}
